% The rule behind each reported number, pinned against bench/runner.py
% and bench/report.py by a test so the list cannot drift from the code
% that produced the tables.

A benchmark number without the rule that produced it is not comparable to
anyone else's, and several of ours have a plausible-looking alternative
reading. The rules below are the ones the harness applies. They are pinned
against \texttt{bench/runner.py} and \texttt{bench/report.py} by a test, so
this list cannot drift from the code that produced the tables.

\paragraph{Everything is scored on true deltas.} In the noisy suites the
measurement itself lies, and the ledger acts on what it was told. Every
metric below is computed from the resource movement that actually occurred;
what the ledger was told is carried separately as
\texttt{reported\_cum\_delta}. Scoring on the reported number would credit an
arm for being deceived in a convenient direction, which is the opposite of
the question those suites ask.

\paragraph{Kill rate} is the fraction of seeds with a kill cycle, and a kill
cycle is one after which no surviving entry from the poisoned source still
advises an action. It is a claim about advice that can be acted on, not about
text remaining in the store.

\paragraph{Poison alive} counts surviving entries carrying the poisoned
source in their provenance, inert ones included. It is reported beside the
kill rate and never instead of it. A mechanism that removes nothing inert
reads a kill rate of $1.00$ and looks immune, and three readings in this
paper were flattered exactly that way before this count existed
(\S\ref{sec:limitations}).

\paragraph{Kill cycle} is the median over the seeds that killed, with a
bootstrap confidence interval; seeds that never killed are excluded from it
rather than entered as a large number. The median therefore describes only
the seeds that got there, and must be read with the kill rate beside it: an
arm that kills once, quickly, reports a better kill cycle than an arm that
kills every time.

\paragraph{Damage before kill} sums the negative true deltas from the start
of the run up to and including the kill cycle, and over the whole run when
there was no kill. It prices what the poison cost before selection removed
it, which is the cost the mechanism is supposed to bound rather than
eliminate.

\paragraph{Cumulative and tail $\Delta$} are the sum of true deltas over the
whole run, and the mean true delta over its last five cycles. The tail is the
steady state after selection has settled; the cumulative figure includes the
price of getting there.

\paragraph{Final population} is the number of living entries when the run
ends. It is a leanness measure, not a safety one: two arms with identical
outcomes and different populations differ in what they cost to carry, not in
what they got right.
